\documentclass[11pt]{article}
\usepackage[margin=1in]{geometry}
\usepackage{amsmath,amssymb,amsfonts}
\usepackage{graphicx}
\usepackage{booktabs}
\usepackage{multirow}
\usepackage{xcolor}
\usepackage{microtype}
\usepackage{etoolbox}
\AtBeginEnvironment{tabular}{\footnotesize}
\renewcommand{\arraystretch}{0.92}
\setlength{\tabcolsep}{4.5pt}
\usepackage{natbib}
\usepackage{xr}
\externaldocument{main}
\usepackage{hyperref}
\title{Supplementary Material\\ \large Risk-Aware Corridor-Constrained Pathfinding for UAV Navigation in Biosecurity-Sensitive Environments}
\author{}\date{}
\begin{document}\maketitle
\noindent This document contains (i)~secondary baseline, benchmark, and feasibility experiments and (ii)~full proofs of the formal results stated in the main paper. Cross-references such as ``Proposition~1'', ``Eq.~(3)'', ``Table~2'', and ``Section~6'' refer to the main paper.

\section{Secondary Experiments}
\subsection*{Experimental setup}

\subsection{Experiment 6: Weighted A* Comparison}

\paragraph{Setup.}
We evaluated Weighted A* (wA*, $w \in \{1.5,\, 2.0,\, 3.0\}$) against A*, ILS, and AILS on $200 \times 200$ risk-annotated grids, using the Experiment~1 density--risk--$\lambda$ configurations with 50 maps each and random start--goal pairs (minimum Manhattan distance 100 cells).

\subsection{Experiment 10: Moving AI Lab Benchmark Evaluation}

\paragraph{Setup.}
For external validity and comparison with published results, we evaluated A*, ILS, and AILS on Moving AI Lab maps~\citep{Sturtevant2012} in three categories: \emph{game maps} from Dragon Age: Origins (lak303d, den501d, brc202d, ost003d); \emph{mazes} (maze512-$k$-0, $k \in \{1, 2, 4, 8\}$, progressively wider corridors); and \emph{random maps} (random512-$d$-0, $d \in \{10, 20, 30, 40\}$\% density) matching the procedural grids of Experiments~1--9. All maps are $512 \times 512$ or smaller; up to 100 scenarios per map were drawn from the official files, preferring longer paths. No risk ($\lambda = 0$): this tests the corridor mechanism on diverse topologies. ILS used $\alpha = 0.05$; AILS used $r_{\min} = 0.02 \cdot \text{diag}$, $r_{\max} = 0.10 \cdot \text{diag}$.

\subsection{Experiment 10b: Risk-Annotated Benchmark Evaluation}

\paragraph{Setup.}
To extend benchmarks beyond $\lambda = 0$, we overlaid synthetic risk (gradient, hotspot) on Moving AI Lab random maps and evaluated ILS, AILS, and RILS at $\lambda \in \{0.5, 1.0\}$, testing two maps (random512-20-0, random512-30-0) with 10 longer-path scenarios per configuration.

\subsection{Experiment 10c: Maze Benchmark Evaluation}

\paragraph{Setup.}
To stress-test corridor methods on adversarial topologies, we evaluated ILS and AILS on procedural $256 \times 256$ maze maps at four corridor widths ($k \in \{1, 2, 4, 8\}$; densities 51\%, 34\%, 21\%, 14\%), 20 scenarios per configuration, no risk ($\lambda = 0$). Mazes were generated via randomised DFS (perfect, fully reachable). Official Moving AI Lab maze maps were unavailable due to persistent HTTP 500 download errors across multiple dates.

\subsection{Experiment 11: Multi-Heuristic Comparison}

\paragraph{Setup.}
To assess robustness to heuristic choice, we compared four heuristics within ILS at $\lambda = 1.0$: octile, Euclidean, Manhattan, and risk-adjusted octile ($h_\lambda = h \cdot (1 + \lambda \cdot 0.3)$), each with A* and ILS ($\alpha = 0.05$) on $200 \times 200$ risk-annotated grids at three densities and three risk distributions (20 maps per configuration).

\subsection{Experiment 13: Autopilot Waypoint-Mission Feasibility}\label{subsec:exp13_setup}

\paragraph{Setup.}
To test whether corridor-restricted paths convert into a deployable autopilot mission, we implemented an ArduPilot-compatible pipeline: grid paths become GPS waypoints (1\,m per cell), compressed by collinear-waypoint removal, and assembled into a mission with a 2\,m proximity threshold and return-to-launch terminator. We evaluate it in a deterministic mission simulator reproducing the waypoint-following geometry; we do \emph{not} claim closed-loop flight in the ArduCopter SITL simulator, the immediate next step (Section~\ref{subsec:practical}). A representative mission used a $50 \times 50$ grid at $20\%$ density.

\paragraph{Parameter sweep.}
We then ran a feasibility sweep (100 maps per configuration) over grid size ($30$--$100$ cells), density ($10\%$--$30\%$), and initial corridor width ($3\%$--$15\%$). Metrics: planning time, node reduction, path optimality ratio, post-compression waypoint count, and mean heading change between consecutive waypoints (smoothness). This covers the small-grid flight-envelope regime, below the AILS wall-clock crossover (${\sim}300 \times 300$) of Experiment~9.


\subsection*{Results}

\subsection{Experiment 6: Weighted A* Comparison}

Table~\ref{tab:exp6_wa} compares Weighted A* (wA*) against ILS across risk distributions at $\lambda = 1.0$, averaged over three densities with 50 random start--goal pairs per configuration.

\begin{table}[!htbp]
\centering
\caption{Weighted A* vs.\ ILS at $\lambda = 1.0$ (mean across three densities, 50 maps each, random start--goal pairs). Speedup relative to A*; Opt = mean path cost ratio $\pm$ 95\% CI.}
\label{tab:exp6_wa}
\resizebox{\textwidth}{!}{%
\begin{tabular}{lcccccccc}
\toprule
\multirow{2}{*}{\textbf{Risk}} & \multicolumn{2}{c}{\textbf{wA* ($w\!=\!1.5$)}} & \multicolumn{2}{c}{\textbf{wA* ($w\!=\!2.0$)}} & \multicolumn{2}{c}{\textbf{wA* ($w\!=\!3.0$)}} & \multicolumn{2}{c}{\textbf{ILS}} \\
& Speed & Opt & Speed & Opt & Speed & Opt & Speed & Opt \\
\midrule
Gradient & $7.2\times$ & $1.008 \pm .001$ & $95.2\times$ & $1.022 \pm .002$ & $90.9\times$ & $1.028 \pm .003$ & $6.6\times$ & $1.000 \pm .000$ \\
Hotspot  & $18.2\times$ & $1.039 \pm .006$ & $22.6\times$ & $1.045 \pm .007$ & $22.0\times$ & $1.051 \pm .007$ & $1.9\times$ & $1.009 \pm .003$ \\
Uniform  & $40.8\times$ & $1.016 \pm .002$ & $48.8\times$ & $1.021 \pm .003$ & $48.6\times$ & $1.026 \pm .003$ & $3.3\times$ & $1.000 \pm .000$ \\
\midrule
Mean     & $22.1\times$ & 1.021 & $55.5\times$ & 1.029 & $53.9\times$ & 1.035 & $3.9\times$ & 1.003 \\
\bottomrule
\end{tabular}%
}
\end{table}

wA* and ILS occupy different points on the speed--quality tradeoff. wA* reached far higher speedups ($22$--$56\times$ vs.\ ILS's $3.9\times$, up to $95\times$ for gradient risk at $w = 2.0$) via heuristic inflation that expanded $92$--$98\%$ fewer nodes, but at $2.0$--$4.6\%$ mean suboptimality (worst on hotspot risk, $3.5$--$4.6\%$). ILS stayed within $0.3\%$ of optimal everywhere, with gradient and uniform risk near-perfect ($\leq 0.01\%$).

For biosecurity, where path quality affects risk exposure, ILS's $6.6\times$ at $0.01\%$ suboptimality on gradient risk is preferable to wA*'s $95.6\times$ at $2.0\%$; for time-critical scenarios tolerating bounded suboptimality, wA* with $w = 1.5$ is an effective compromise ($7.5$--$43\times$, $\leq 3.5\%$ suboptimal).

\subsection{Experiment 10: Moving AI Lab Benchmark Evaluation}

\begin{table}[!htbp]
\centering
\caption{ILS and AILS performance on Moving AI Lab benchmark maps ($\lambda = 0$, up to 100 scenarios per map). Speed = speedup vs.\ A*; Red = node reduction \%; Opt = path cost ratio vs.\ A*; Slv = solve rate.}
\label{tab:exp10_benchmark}
\resizebox{\textwidth}{!}{%
\begin{tabular}{llcccccccc}
\toprule
\multirow{2}{*}{\textbf{Cat.}} & \multirow{2}{*}{\textbf{Map}} & \multicolumn{4}{c}{\textbf{ILS}} & \multicolumn{4}{c}{\textbf{AILS}} \\
& & Speed & Red.\% & Opt & Slv\% & Speed & Red.\% & Opt & Slv\% \\
\midrule
Game & lak303d & $0.29\times$ & $-231$ & 1.000 & 58 & $0.35\times$ & $-168$ & 1.001 & 88 \\
Game & den501d & $0.30\times$ & $-234$ & 1.001 & 94 & $0.37\times$ & $-163$ & 1.001 & 94 \\
Game & brc202d & $0.37\times$ & $-170$ & 1.000 & 49 & $0.47\times$ & $-106$ & 1.000 & 50 \\
Game & ost003d & $0.28\times$ & $-232$ & 1.000 & 50 & $0.32\times$ & $-197$ & 1.001 & 50 \\
\midrule
\multicolumn{2}{l}{\textbf{Game mean}} & $0.31\times$ & $-216$ & 1.000 & 63 & $0.38\times$ & $-159$ & 1.001 & 70 \\
\midrule
Random & 10\% & $0.94\times$ & $0.1$ & 1.000 & 100 & $0.89\times$ & $4.2$ & 1.000 & 100 \\
Random & 20\% & $0.93\times$ & $0.1$ & 1.000 & 100 & $0.89\times$ & $1.1$ & 1.000 & 100 \\
Random & 30\% & $0.95\times$ & $1.3$ & 1.000 & 100 & $0.93\times$ & $1.6$ & 1.000 & 100 \\
Random & 40\% & $1.10\times$ & $12.5$ & 1.000 & 100 & $0.94\times$ & $5.4$ & 1.000 & 100 \\
\midrule
\multicolumn{2}{l}{\textbf{Random mean}} & $0.98\times$ & $3.5$ & 1.000 & 100 & $0.91\times$ & $3.1$ & 1.000 & 100 \\
\bottomrule
\end{tabular}%
}
\end{table}

Table~\ref{tab:exp10_benchmark} shows corridor methods give no speedup on standard benchmarks at $\lambda = 0$ (no risk layer), and are substantially slower on game maps.

\paragraph{Game maps (Dragon Age: Origins).}
On game maps (60--83\% obstacle density), ILS was $2.7$ to $3.6$ times slower than A* ($0.28\times$--$0.37\times$) with strongly negative node reduction ($-170\%$ to $-234\%$); AILS was slightly better ($0.32\times$--$0.47\times$) but still slower. Game maps are rooms joined by narrow doorways, so the Bresenham line crosses many walls and each failed corridor attempt accumulates wasted search beyond A*'s single pass. Solve rates were lower (49 to 94\%) where scenarios require paths far from the Bresenham line.

\paragraph{Random maps (10--40\% density).}
On random maps at $\lambda = 0$, ILS was near-neutral ($0.93$--$1.10\times$), with positive node reduction only at 40\% density ($12.5\%$), consistent with Experiment~1's $0.90$--$1.37\times$: on uniform-cost grids the optimal path lies near the geometric shortest path, leaving little for the corridor to prune. All scenarios solved optimally (cost ratio $\leq 1.0001$).

\paragraph{Implications.}
This is the clearest evidence for a claim that recurs throughout the paper: the corridor advantage derives from the risk-weighted cost landscape, not from the corridor mechanism alone. On uniform-cost benchmarks it adds overhead without proportional benefit; the $3\times$ to $8\times$ speedups of Experiments 1 and 6 to 8 arise because the risk-weighted cost ($\lambda > 0$) makes a spatial corridor prune the nodes A* explores to evaluate risk-avoiding detours. Corridor methods are thus a targeted accelerator for risk-weighted pathfinding, not a general-purpose speedup.

\paragraph{Full-corpus extension and a topological boundary.}
To test this at scale we evaluated the \emph{entire} Dragon Age: Origins corpus --- all $156$ maps, $7{,}800$ scenarios (longest paths per map), suboptimality scored against the published \texttt{.scen} optima under the no-corner-cutting convention. At $\lambda = 0$ the corridor methods are slower than A* as expected. Adding a gradient risk layer ($\lambda = 1.0$) raises their solve rate (ILS $56.0\%$, AILS $50.4\%$) and keeps path quality near-optimal (mean suboptimality $0.09\%$ for ILS, $0.16\%$ for AILS), but the mean node-expansion reduction \emph{remains negative} (ILS $-39.7\%$, AILS $-11.9\%$). This refines the activation claim of Experiment~10b: risk weighting activates the corridor advantage on \emph{open/random} topologies (where the optimal path tracks the start--goal line), but \emph{not} on the room-and-corridor topology of DAO, where the optimal path deviates far from the Bresenham line. The containment analysis of Section~\ref{subsec:containment} explains why: DAO instances have a large minimal deviation $\zeta^\star$, so the corridor has to widen a long way (or give up) before it contains an optimal path, and weighting the cost by risk does nothing to shrink that geometric deviation. The corridor advantage thus requires two preconditions: a risk-weighted cost function and a topology in which the optimal path stays near the reference line.

\subsection{Experiment 10b: Risk-Annotated Benchmark Results}

\begin{table}[!htbp]
\centering
\caption{ILS, AILS, and RILS on Moving AI Lab random maps with synthetic risk layers. Speed = speedup vs.\ A*; Red = node reduction \%; Opt = path cost ratio.}
\label{tab:exp10b_risk_bench}
\resizebox{\textwidth}{!}{%
\begin{tabular}{llccccccccc}
\toprule
\multirow{2}{*}{\textbf{Map}} & \multirow{2}{*}{$\boldsymbol{\lambda}$} & \multicolumn{3}{c}{\textbf{ILS}} & \multicolumn{3}{c}{\textbf{AILS}} & \multicolumn{3}{c}{\textbf{RILS}} \\
& & Speed & Red.\% & Opt & Speed & Red.\% & Opt & Speed & Red.\% & Opt \\
\midrule
\multicolumn{11}{l}{\textit{Gradient risk}} \\
20\% & 0.5 & $1.71\times$ & 41.2 & 1.000 & $2.18\times$ & 51.9 & 1.000 & $1.11\times$ & 11.7 & 1.000 \\
20\% & 1.0 & $2.57\times$ & 56.7 & 1.000 & $3.02\times$ & 65.0 & 1.000 & $1.39\times$ & 26.5 & 1.000 \\
30\% & 0.5 & $1.96\times$ & 46.2 & 1.000 & $1.92\times$ & 46.8 & 1.000 & $1.07\times$ & 11.6 & 1.000 \\
30\% & 1.0 & $2.68\times$ & 59.6 & 1.000 & $2.62\times$ & 60.0 & 1.000 & $1.41\times$ & 29.2 & 1.000 \\
\midrule
\multicolumn{11}{l}{\textit{Hotspot risk}} \\
20\% & 0.5 & $1.42\times$ & 27.0 & 1.000 & $1.75\times$ & 39.6 & 1.000 & $1.19\times$ & 17.5 & 1.000 \\
20\% & 1.0 & $2.01\times$ & 49.4 & 1.000 & $2.49\times$ & 59.1 & 1.001 & $1.56\times$ & 34.2 & 1.000 \\
30\% & 0.5 & $1.28\times$ & 25.4 & 1.000 & $1.36\times$ & 26.8 & 1.000 & $1.11\times$ & 11.8 & 1.000 \\
30\% & 1.0 & $1.83\times$ & 41.1 & 1.000 & $1.86\times$ & 42.5 & 1.000 & $1.29\times$ & 19.3 & 1.000 \\
\bottomrule
\end{tabular}%
}
\end{table}

Table~\ref{tab:exp10b_risk_bench} shows the same Moving AI Lab maps that gave no benefit at $\lambda = 0$ yield substantial speedups once risk layers are added. At $\lambda = 1.0$ on gradient risk, ILS reached $2.57\times$ to $2.68\times$ ($57\%$ to $60\%$ node reduction) and AILS $2.62\times$ to $3.02\times$; on hotspot risk, $1.83\times$ to $2.01\times$ (ILS) and $1.86\times$ to $2.49\times$ (AILS). Path quality stays near-optimal (cost ratio $\leq 1.001$), and RILS holds perfect optimality at lower speedup ($1.07\times$ to $1.56\times$). Results are reproducible: node reduction had zero variance across the $n = 10$ scenarios per configuration, and paired mean speedups across the two maps are close. This isolates the risk-weighted cost landscape as the trigger for corridor acceleration, independent of map topology.

\subsection{Experiment 10c: Maze Benchmark Results}

\begin{table}[!htbp]
\centering
\caption{ILS and AILS performance on procedurally generated maze maps ($256 \times 256$, $\lambda = 0$). Solve rate denotes the fraction of scenarios where the corridor method found a path.}
\label{tab:exp10c_maze}
\resizebox{\textwidth}{!}{%
\begin{tabular}{lcccccccccc}
\toprule
\multirow{2}{*}{\textbf{Maze}} & \multirow{2}{*}{\textbf{Width}} & \multirow{2}{*}{\textbf{Dens.}} & \multicolumn{4}{c}{\textbf{ILS}} & \multicolumn{4}{c}{\textbf{AILS}} \\
& & & Speed & Red.\% & Opt & Slv\% & Speed & Red.\% & Opt & Slv\% \\
\midrule
maze-1 & 1 & 51\% & $0.47\times$ & $-105$ & 1.000 & 45 & $0.43\times$ & $-123$ & 1.000 & 50 \\
maze-2 & 2 & 34\% & $0.65\times$ & $-50$ & 1.000 & 45 & $0.63\times$ & $-54$ & 1.000 & 45 \\
maze-4 & 4 & 21\% & $0.53\times$ & $-74$ & 1.000 & 40 & $0.54\times$ & $-64$ & 1.000 & 35 \\
maze-8 & 8 & 14\% & $0.43\times$ & $-134$ & 1.000 & 65 & $0.46\times$ & $-111$ & 1.000 & 55 \\
\midrule
\multicolumn{3}{l}{\textbf{Mean}} & $0.52\times$ & $-91$ & 1.000 & 49 & $0.51\times$ & $-88$ & 1.000 & 46 \\
\bottomrule
\end{tabular}%
}
\end{table}

Mazes (Table~\ref{tab:exp10c_maze}) are a worst case: both ILS and AILS were slower than A* ($0.43$--$0.67\times$) with massively negative node reduction ($-50\%$ to $-134\%$), since maze paths wind through narrow passages far from the Bresenham line and each failed corridor attempt wastes search before expanding to contain the tortuous optimum. Solve rates were low ($35$--$65\%$), as many paths cannot be found within the maximum 10 corridor expansions. Path quality nonetheless stayed perfect ($1.000$) for every solved scenario. This delimits the applicability of the methods: corridor search is effective where the Bresenham line approximates the optimal path, and fails on topologies such as mazes where the optimum bears no geometric relation to the start--goal line.

\subsection{Experiment 11: Multi-Heuristic Results}

\begin{table}[!htbp]
\centering
\caption{ILS speedup with different heuristics at $\lambda = 1.0$ on $200 \times 200$ risk-annotated grids (mean across three densities).}
\label{tab:exp11_heuristic}
\begin{tabular}{lcccccc}
\toprule
\multirow{2}{*}{\textbf{Heuristic}} & \multicolumn{3}{c}{\textbf{ILS Speedup ($\times$)}} & \multicolumn{3}{c}{\textbf{ILS Node Red.\ (\%)}} \\
& Grad. & Hot. & Unif. & Grad. & Hot. & Unif. \\
\midrule
Octile         & $1.91$ & $1.60$ & $1.57$ & $48.9$ & $36.4$ & $36.3$ \\
Euclidean      & $2.11$ & $1.64$ & $1.66$ & $52.2$ & $39.4$ & $41.2$ \\
Manhattan      & $1.63$ & $1.35$ & $1.28$ & $35.0$ & $29.8$ & $25.2$ \\
Risk-adj.      & $1.23$ & $1.15$ & $1.02$ & $24.8$ & $15.0$ & $6.8$ \\
\bottomrule
\end{tabular}
\end{table}

Table~\ref{tab:exp11_heuristic} shows corridor speedups are robust across heuristics. All four gave positive ILS speedups on gradient and hotspot risk; Euclidean was highest ($2.11\times$ gradient, $1.66\times$ uniform) and Manhattan moderate ($1.28$--$1.63\times$), its coarser guidance exploring more nodes and leaving less to prune. The risk-adjusted heuristic ($h_\lambda = h \cdot (1 + \lambda \cdot 0.3)$) gave the lowest speedups despite guiding A* more efficiently: a tighter heuristic already prunes well, leaving less for the corridor to eliminate, so the advantage is largest when the baseline search is expensive. On uniform risk it reached only $1.02\times$, A* expanding just $312$--$843$ nodes so corridor overhead nearly cancels the benefit. The corridor thus interacts multiplicatively with heuristic quality, giving consistent relative improvement across all heuristics; the default octile heuristic is a practical compromise.

\subsection{Experiment 13: Autopilot Waypoint-Mission Feasibility Results}\label{subsec:exp13_results}

The $50 \times 50$ ILS path converted cleanly into a deployable mission: a compact GPS waypoint sequence with a 2\,m proximity threshold and return-to-launch terminator, with smooth heading transitions and no post-processing beyond collinear-waypoint compression. A 100-map sweep over grid size ($30$--$100$), density ($10$--$30\%$) and corridor width ($3$--$15\%$) confirmed the missions stay geometrically clean: optimality within $0.2\%$ of A* up to $25\%$ density (cost ratio $1.0001$ at $10\%$ to $1.0019$ at $25\%$), degrading mildly to $1.0035$ at $30\%$; collinear compression removed $\approx 69\%$ of waypoints, with mean heading change $\approx 13^\circ$. Compression and smoothness were independent of corridor width, while optimality improved monotonically with it ($1.0035$ at a $3\%$ band to $1.0000$ at $15\%$). On these sub-$100 \times 100$ grids planning time was only modestly below A* (mean $1.14\times$, range $0.87\times$--$2.19\times$); the substantial gains of Experiments~1, 9 and~14 emerge above the $\sim 300 \times 300$ crossover. Closed-loop SITL flight is the next step (Section~\ref{subsec:practical}).


\section{Proofs}\label{app:proofs-supp}
Proofs are given in the order the results are stated in the main paper.


\paragraph{Proposition~\ref{prop:admissible}.}
\noindent\textit{Proof.}
For any path $\pi = (n, v_1, \ldots, g)$ with $k$ edges, the true cost decomposes as
\[
C(\pi) = \sum_{i=0}^{k-1} \bigl[d(v_i, v_{i+1}) + \lambda \cdot \rho(v_{i+1})\bigr] = D(\pi) + \lambda \cdot R(\pi),
\]
where $D(\pi) \geq h(n,g)$ because the octile distance is the minimum-cost path on the unblocked 8-connected grid, and $R(\pi) = \sum_{i=1}^{k} \rho(v_i) = k \cdot \bar{\rho}_\pi$ by definition of the mean risk.
Under condition~(\ref{eq:adm_condition}), $\lambda \cdot R(\pi) \geq \lambda \cdot h(n,g) \cdot \bar{\rho}_0$, so
\[
C(\pi) \;=\; D(\pi) + \lambda \cdot R(\pi) \;\geq\; h(n,g) + \lambda \cdot h(n,g) \cdot \bar{\rho}_0 \;=\; h_\lambda(n,g).
\]
Therefore $h_\lambda$ is admissible at $n$. $\square$

\paragraph{Proposition~\ref{prop:complete}.}
\noindent\textit{Proof.}
When $\pi^* \subseteq \mathcal{C}(r)$, the corridor restriction removes only nodes not on $\pi^*$, and A* with an admissible heuristic on the restricted graph $G|_{\mathcal{C}(r)}$ finds $\pi^*$ because all nodes on $\pi^*$ remain in the search space.
For completeness, each failed attempt increases corridor width by $\Delta r$, so after $K_{\max}$ attempts the corridor covers all cells within $d_\infty$-distance $r_0 + (K_{\max}-1) \cdot \Delta r$ of the reference line.
$\square$

\paragraph{Proposition~\ref{prop:termination}.}
\noindent\textit{Proof.}
The half-width updates monotonically as $r \leftarrow r + \Delta r$ on each failed attempt, so at attempt $k$ the corridor is $\mathcal{C}(r_0 + (k-1)\Delta r)$; the loop therefore stops after at most $\min\!\big(K_{\max},\, \lceil (r_{\mathrm{cover}} - r_0)/\Delta r \rceil + 1\big)$ iterations, the second term being the number of widenings needed to reach $r_{\mathrm{cover}}$. Each restricted search terminates because the base algorithm is complete on the finite subgraph $G|_{\mathcal{C}(r)}$ and returns its in-corridor optimum or reports infeasibility, so on exhaustion the procedure returns the optimum of the terminal corridor $\mathcal{C}(r_0+(K_{\max}-1)\Delta r)$ or FAILURE. For the completeness claim, every cell of an $H \times W$ grid lies within $d_\infty$-distance $\max(H,W) - 1$ of either endpoint of the reference line $\ell(s,g)$, hence within that distance of $\ell(s,g)$; so when the cap permits $r$ to reach $r_{\mathrm{cover}} = \max(H,W) - 1$ the corridor covers $V$ (the choice $\min(H,W)$ is insufficient on non-square grids whose reference line is short and corner-localised), and that final attempt runs the base algorithm on the unrestricted graph $G$, returning a path iff one exists. If the cap halts the loop before $r_{\mathrm{cover}}$, the same per-attempt argument yields the terminal-corridor optimum, or FAILURE if it contains no $s$-to-$g$ path. The pathological case $\Delta r = 0$ (a corridor that never widens) is excluded by the hypothesis $\Delta r \geq 1$, which is the natural integer increment on an 8-connected grid with unit spacing. $\square$

\paragraph{Proposition~\ref{prop:certificate}.}
\noindent\textit{Proof.}
Optimality of $\pi_{\mathrm{ILS}}$ on $G|_{\mathcal{C}^{\dagger}}$ is inherited from the base algorithm's optimality on any finite weighted graph. Its cost therefore equals the minimum in-corridor cost, which exceeds the global optimum $C(\pi^*)$ by exactly $\Delta_{\mathcal{C}^{\dagger}}$; this difference is zero precisely when a globally optimal path lies inside $\mathcal{C}^{\dagger}$, and trivially when $r^{\dagger} = r_{\mathrm{cover}}$. Inequality~\eqref{eq:cert_loose} follows from $C(\pi^*) \geq \underline{C}(s,g)$; equality~\eqref{eq:cert_exact} follows by computing $C(\pi^*)$ exactly with one unrestricted search. $\square$

\paragraph{Corollary~\ref{cor:bound}.}
\noindent\textit{Proof.}
The base algorithm's optimality on $G|_{\mathcal{C}^{\dagger}}$ (as in Proposition~\ref{prop:certificate}) makes $\pi_r^* = \pi_{\mathrm{ILS}}$ the minimum-cost in-corridor path; it uses at most $\ell_{\max}(\mathcal{C}^{\dagger})$ edges, each of weight at most $w_{\max}(G)$, and $C(\pi^*) \geq 1 > 0$ since every edge weight is at least $1$ (so the division in Eq.~\eqref{eq:additive_bound} is well defined). $\square$

\paragraph{Theorem~\ref{thm:no_bound}.}
\noindent\textit{Proof (explicit witness family, respecting $w = d + \lambda\rho \geq 1$).}
Fix $\epsilon > 0$, choose risk weight $\lambda > \epsilon$, and set $K := r_0 + 2$. Build a $(K{+}2)\times N$ grid (rows $0,\dots,K{+}1$; columns $0,\dots,N{-}1$) in the paper's $8$-connected model, with $w(u,v) = d(u,v) + \lambda\rho(v)$, $\rho \in [0,1]$: row~$1$ is free with \emph{maximal risk} $\rho \equiv 1$; row~$0$ is blocked; rows $2,\dots,K{-}1$ are blocked except at the two end columns $0$ and $N{-}1$, whose cells are free with risk~$0$ (the only vertical connectors); row~$K$ is the free \emph{safe lane} with risk~$0$; $s = (1,0)$, $g = (1,N{-}1)$. With row~$0$ blocked and rows $2,\dots,K{-}1$ blocked except at the ends, the diagonal moves available in the $8$-connected model create no in-corridor shortcut: from a row-$1$ cell every diagonal lands either in blocked row~$0$ or in a blocked interior cell of row~$2$, so in-corridor travel is confined to horizontal steps along row~$1$ (each costing $1+\lambda$). All edge costs lie in $[1, \sqrt2 + \lambda]$, a legal instance. With $r_0 = K{-}2$ the corridor $\mathcal{C}(r_0)$ covers rows $1,\dots,K{-}1$ but not row~$K$ (Chebyshev distance $K{-}1 = r_0{+}1$ from row~$1$); the in-corridor end-column connectors are dead-ends (they reach only the excluded row~$K$), so the only in-corridor $s$-to-$g$ path runs along row~$1$, of cost $(1+\lambda)(N{-}1) = \Theta(N)$. A path exists at $r_0$, so ILS returns it without widening: $C(\pi_{\mathrm{ILS}}) = (1+\lambda)(N{-}1)$. The global optimum climbs the column-$0$ connector to the safe lane (cost $K{-}1$), traverses row~$K$ (cost $N{-}1$), and descends the column-$(N{-}1)$ connector to $g$ (cost $K{-}1$, the final step into the risk-$1$ cell $g$ costing $1+\lambda$): $C(\pi^*) \leq (N{-}1) + 2(K{-}1) + \lambda$. Hence
\[
\frac{C(\pi_{\mathrm{ILS}})}{C(\pi^*)} \;\geq\; \frac{(1+\lambda)(N{-}1)}{(N{-}1) + 2(K{-}1) + \lambda} \;\longrightarrow\; 1 + \lambda \quad (N \to \infty),
\]
which exceeds $1+\epsilon$ for $N$ large (the safe-lane detour is the global optimum once $\lambda(N{-}1) > 2(K{-}1) + \lambda$, i.e.\ for $N$ large). The construction is in the spirit of the no-uniform-suboptimality argument of \citet{Pohl1970}, transplanted to the geometric corridor-restriction setting. The two paths compared here use only cardinal moves with per-step cost in $[1, 1+\lambda]$, so for \emph{fixed} $\lambda$ the cost ratio is bounded by $1+\lambda$; unboundedness over the family therefore comes from $\lambda \to \infty$, not from $N$, which is the precise sense in which no risk-weight-independent $(1+\epsilon)$ bound holds. A direct cell-count check ($K=3$, $\lambda=9$, $N=1000$) gives $\frac{10 \cdot 999}{999 + 4 + 9} = \frac{9990}{1012} = 9.87$, consistent with the asymptotic $1+\lambda = 10$. $\square$

\paragraph{Theorem~\ref{thm:size}.}
\noindent\textit{Proof.}
Assume w.l.o.g.\ the dominant axis is the column axis, so the Bresenham line has one cell per column $c \in \{c_{\min},\dots,c_{\max}\}$ with $c_{\max} - c_{\min} = D_\infty$ and row value $y(c)$ monotone and $1$-Lipschitz. Every corridor cell lies in one of at most $D_\infty + 2r + 1$ columns. Fix a column $j$: a cell $(j,y) \in \mathcal{C}(r)$ requires a line cell $(c,y(c))$ with $c \in [j-r,j+r] \cap [c_{\min},c_{\max}] =: [a,b]$ and $y \in [y(c)-r, y(c)+r]$. By monotonicity and $1$-Lipschitzness, $\{y(c): c \in [a,b]\} \subseteq [y(a), y(b)]$ with $y(b) - y(a) \leq b - a \leq 2r$, so the covered rows lie in $[y(a)-r, y(b)+r]$, at most $4r+1$ integers. Multiplying by the column count gives the bound; for an axis-aligned line $y(c)$ is constant and the per-column count is $2r+1$. $\square$

\paragraph{Corollary~\ref{cor:reduction}.}
\noindent\textit{Proof.}
A consistent-heuristic A* (or Dijkstra) settles each vertex once, so expansions $\leq |\mathcal{C}(r)|$; the heap performs $O(|\mathcal{C}(r)|)$ operations each $O(\log|\mathcal{C}(r)|)$. Theorem~\ref{thm:size} gives $|\mathcal{C}(r)| = O(\alpha N^2)$, and the matching lower bound $|\mathcal{C}(r)| \geq (2r+1)\,D_\infty = \Omega(\alpha N^2)$ holds because the $\geq 2r+1$ cells stacked on each of the $D_\infty+1$ reference-line columns are distinct. Hence $|\mathcal{C}(r)| = \Theta(\alpha N^2)$, and the reduction relative to the $\Theta(N^2)$ worst case is $\Theta(1/\alpha)$; with $h_\lambda$, which is not consistent, the per-vertex single-expansion bound requires $h_0$ or Dijkstra. $\square$

\paragraph{Theorem~\ref{thm:complexity}.}
\noindent\textit{Proof.}
(i) Bresenham emits one cell per dominant-axis step; the summed-area table is a two-pass prefix sum over $HW$ entries, each rectangular query four look-ups. (ii) A dense mask initialises $HW$ entries; the base search runs on $|\mathcal{C}(r_k)| \leq HW$ vertices. (iii) The half-width increases by $\Delta r \geq 1$ per failed attempt and caps at $K_{\max}$; every cell is within Chebyshev distance $N_{\!*}-1$ of an endpoint of $\ell(s,g)$, hence of $\ell(s,g)$, so $\mathcal{C}(N_{\!*}-1) \supseteq V$. (iv) Sum (ii) over $O(N_{\!*})$ iterations with $|\mathcal{C}(r_k)| \leq HW$; one mask and one table are reused, giving $O(HW)$ space. $\square$

\paragraph{Lemma~\ref{lem:containment}.}
\noindent\textit{Proof.}
Let $\pi^*$ attain $\zeta(\pi^*) = \zeta^\star \leq r$; then $\pi^* \subseteq \mathcal{C}(r)$ and its edges lie in $G|_{\mathcal{C}(r)}$. The base search is optimal on the finite graph $G|_{\mathcal{C}(r)}$, so it returns a path of cost $\leq C(\pi^*)$, hence equal to the global minimum. $\square$

\paragraph{Corollary~\ref{cor:freeline}.}
\noindent\textit{Proof.}
For $\lambda = 0$, $C(\pi) \geq h(s,g)$ for every path. The Bresenham line uses $\min(|\Delta x|,|\Delta y|)$ diagonal and $\bigl||\Delta x|-|\Delta y|\bigr|$ cardinal steps, of total cost $\sqrt2 \min(|\Delta x|,|\Delta y|) + \bigl||\Delta x|-|\Delta y|\bigr| = h(s,g)$; being obstacle-free it is a geodesic in $V$, so $\zeta^\star = 0$. $\square$

\paragraph{Proposition~\ref{prop:isolated_feas}.}
\noindent\textit{Proof.}
The path that follows $\ell$ until meeting $B$, routes along the outer ring $\{v: d_\infty(v,p_0) = b+1\}$, then rejoins $\ell$, has deviation $\leq b+1$ (ring cells satisfy $d_\infty(v,\ell) \leq d_\infty(v,p_0) = b+1$) and is feasible, so $G|_{\mathcal{C}(b+1)}$ admits an $s$-to-$g$ path and the base search returns the in-corridor optimum. $\square$

\paragraph{Lemma~\ref{lem:jps_unsound}.}
\noindent\textit{Proof.}
Take the obstacle-free $3\times4$ grid, $s = (2,0)$, $g = (0,3)$, $\lambda = 1$, $\rho(0,2) = 1$ and $\rho \equiv 0$ elsewhere. The octile distance $h(s,g) = 2\sqrt2 + 1$ is realised by exactly three monotone geodesics; two of them, $P_2 = (2,0)(1,1)(1,2)(0,3)$ and $P_3 = (2,0)(2,1)(1,2)(0,3)$, avoid $(0,2)$ and cost $2\sqrt2+1$, while $P_1 = (2,0)(1,1)(0,2)(0,3)$ costs $2\sqrt2 + 2$. Since the grid is obstacle-free, no cell has a forced neighbour, so a straight scan halts only at the goal or a boundary and a diagonal scan generates a successor only where a straight sub-scan finds a jump point or the goal is intercepted. We trace every jump-point successor reachable from $s=(2,0)$ toward $g$. The cardinal scans from $s$ dead-end: rightward along row~$2$ reaches the boundary with no forced neighbour (so $(2,1)$ is scanned but \emph{not} generated as a successor and is never itself expanded), and upward along column~$0$ likewise. The diagonal (up-right) scan from $s$ steps to $(1,1)$, where its straight sub-scans (rightward through $(1,2),(1,3)$; upward through $(0,1)$) find no forced neighbour and no goal interception, so it continues to $(0,2)$ on the top boundary; from $(0,2)$ the rightward scan intercepts the goal column and yields $g=(0,3)$. Thus the only jump-point successors inserted are $(0,2)$ and $g$ (the diagonal scan passes through $(1,1)$ without inserting it), neither of which is $(1,2)$; and because $(2,1)$ is never expanded, no diagonal from $(2,1)$ to $(1,2)$ is ever launched. Hence $(1,2)$ is never generated as a jump-point successor. Both optimal paths $P_2,P_3$ pass through $(1,2)$, so neither lies in the jump-point graph, and every $s$-to-$g$ route in that graph passes through $(0,2)$ at cost $\geq 2\sqrt2 + 2 > 2\sqrt2 + 1$. JPS therefore returns a strictly suboptimal path. $\square$

\begin{thebibliography}{9}
\bibitem[Pohl(1970)]{Pohl1970}
Pohl, I., 1970.
Heuristic search viewed as path finding in a graph.
Artificial Intelligence 1(3--4), 193--204.

\bibitem[Sturtevant(2012)]{Sturtevant2012}
Sturtevant, N.R., 2012.
Benchmarks for grid-based pathfinding.
IEEE Transactions on Computational Intelligence and AI in Games 4(2), 144--148.
\end{thebibliography}
\end{document}
